\chapter{Configuration}
\label{chap:configuration}

Configuration expresses the workload and capacity a node is intended to support.
A security limit needs an owner, a unit and an admission rule, but its chosen
capacity also needs an explanation and a way to change it. Resource ceilings
are therefore collected in a typed kernel policy, separate from the allocation
and cleanup code that enforces them.

This chapter describes current configuration mechanisms and how to qualify a
change. The matching operator guide is
\repofile|docs/guides/configuration.md|; the complete resource-policy inventory
and dimensions are in \repofile|docs/reference/resource-policy.md|.

\section{Configuration sources and authority}

The following inputs have different authorities and lifetimes:
\begin{description}
    \item[Kernel resource policy.] Trusted Rust build input selects capacity
        ceilings, ordinary shares and platform reserves. It is compiled into
        the kernel and remains fixed throughout that boot.
    \item[Bootstrap assignments.] The kernel build consumes
        \path{CATTEN_CLUSTER_SERVICES}. Boot runners accept repeatable
        \texttt{-{}-cluster-service [NAME=]IPv4:PORT} assignments. The first
        committed signed ingress policy replaces the bootstrap assignments.
    \item[Trust policy.] The kernel's one prepared boot-trust policy supplies
        launch and grant verification. Capacity configuration cannot install
        keys, grant platform classification or enable production trust.
        Production remains disabled pending protected provisioning and custody.
    \item[Application launch policy.] Signed deployment and kernel admission
        determine per-domain stack/thread limits, executable variants,
        delegated capabilities and protocol-specific launch profiles. An
        application's request does not enlarge a node-wide capacity pool.
    \item[Persistent state.] A storage format tag and committed state are not
        disposable configuration. Changing capacity never authorizes silent
        reformatting, resetting election state or discarding retained charges.
\end{description}

Internal protocol fields, architectural address ranges, page granularity and
hardware register widths are separate structural contracts. Changing their
numbers requires a coordinated implementation and format review. They are not
ordinary operator capacity settings.

\section{Selecting the compile-time resource policy}

The default settings are in
\repofile|crates/catten/src/resource_policy/defaults.rs|. They retain the existing
capacity values. Operators may edit this central file or select an external
profile with \texttt{CATTEN\_RESOURCE\_POLICY}, which must name an existing file
by absolute path. Unset selects the default; empty, relative, missing and
directory paths reject. The file is trusted Rust source and belongs in the
reviewed build inputs, not in an application manifest or upload request.

A profile defines \texttt{POLICY: Policy} in the policy module's scope. It can
start from \texttt{DEFAULT} and override related fields. The checked-in
\repofile|crates/catten/src/resource_policy/larger.rs| demonstrates co-scaling
scratch-node, snapshot-buffer and mapping-node byte/count capacities by one
quarter while retaining their reserve proportions.

\begin{lstlisting}[language=bash]
CATTEN_RESOURCE_POLICY="$PWD/ops/resource-policy.rs" \
  scripts/run-aarch64.sh --security-test --http-test
\end{lstlisting}

The runners pass this environment input through to the kernel build. Cargo
tracks the selected path, its contents and environment selection. Changing or
unsetting the selection rebuilds the kernel; userspace service images do not
need rebuilding for these policy fields. Boot serial output identifies the
selected policy name. Record the exact profile source and kernel artifact
hashes with qualification evidence; the name alone is not an artifact identity.

There is currently no resource-policy boot-module parser, capability, syscall
or live-resizing interface. The compiled policy does not change after charges
exist. Boot-time resource input is a future installation mechanism, not an
implemented fallback for an invalid build profile.

\section{Related limits and tuning consequences}

The policy groups domain/ordinary/total capacities and resource dimensions.
For fixed pools, every ordinary dimension is positive and below the total;
domain dimensions must fit the ordinary pool. A reserve is the difference
between total and ordinary admission, not a guaranteed share for one service.

Examples of settings that require joint review include:
\begin{itemize}
    \item Memory-object size, aggregate domain/principal backing and snapshot
        bytes. One maximum-sized object must fit an ordinary frame snapshot.
        Increasing object size can also increase mapping work and pin lifetime.
    \item Mapping nodes, scratch extents and snapshots. Concurrent preparations,
        live objects and uncertain cleanup all consume admission. Raising the
        live-object target without space for preparation can prevent progress.
    \item Endpoint counts, queue slots, completion queues, connections,
        pending/reply records, capabilities, waiter entries and observer lists.
        One outstanding RPC can consume several of these families together.
        Endpoint backing rounds queue capacity to a power of two.
    \item Lifecycle watches, publication counts and replacement overlap.
        A name-service replacement needs the new watch before disposing the old
        one; service minima must be tested as part of boot qualification.
    \item Unified memory authority and shared payload/namespace/read-borrower
        node bytes/counts. Every overlapping preparation needs two temporary
        nodes, or three for a read loan, including the namespace candidate for
        an already-established namespace. Published and
        retained metadata consumes its original quota until deallocation.
    \item Thread bundles, principal thread metadata, stacks and private tables.
        More runnable domains can encounter any of these ceilings before their
        individual thread limits, including while old generations remain retained.
    \item Requested metadata bytes, allocator overhead, kernel-heap backing and
        physical progress reserves. Independent logical pool totals do not form
        a complete physical-RAM ledger and cannot establish fairness by addition.
\end{itemize}

RAM-derived pools use documented ratios of discovered usable pages, with
checked widened arithmetic. The heap reserve has a RAM divisor, aligned floor
and aligned ceiling. It clamps before alignment rounding, including extreme
reported RAM. The current launch ABI still bounds heap virtual capacity; the
kernel checks policy against that contract and the actual allocator arena.

Capacity changes do not alter sponsorship or cleanup. A move keeps its original
charge, and root retirement, ASID reuse or later platform promotion cannot refund
live backing. Failed physical cleanup retains the original complete owner and
charge. A larger budget never supplies proof of completion or permission to
replay accepted work.

\section{Planning and validating a policy change}

Use the following process before selecting a profile for a new workload:
\begin{enumerate}
    \item State the workload: RAM, domains/principals, threads, service topology,
        maximum object sizes, outstanding RPCs and expected overlap during restart.
    \item Inventory each consumed family, unit, sponsor, lifetime and rejection
        result. Include unused candidates, detached work, weak references and
        retained uncertain cleanup. Separate structural bounds from capacity policy.
    \item Measure peak occupancy and rejection under the intended concurrency.
        Choose ordinary capacity and trusted headroom together, explaining
        dependencies and expected failure behavior. A previously passing default
        is not a capacity estimate for a larger workload.
    \item Run the policy harness and strict kernel builds. The harness compiles
        the actual policy, exercises admission counters and rejects inconsistent
        profiles. Compile-time validation checks relationships and arithmetic;
        it does not establish boot viability or essential-service progress.
    \item Boot the exact selected profile on the supported backends. Exercise
        success, exhaustion, allocation rejection, cancellation, restart/generation
        reuse, retained charges and confirmed cleanup. Measure management,
        timer/scheduler and teardown progress against deadlines fixed beforehand.
    \item Preserve the profile, artifact hashes, occupancy/latency results and
        failure diagnostics. Keep failed or interrupted attempts separate from
        successful repeats. Requalify after changing the workload or profile.
\end{enumerate}

\begin{lstlisting}[language=bash]
CATTEN_RESOURCE_POLICY="$PWD/ops/resource-policy.rs" \
  python3 scripts/tests/test-resource-policy.py
cargo fmt --all -- --check
\end{lstlisting}

The policy harness checks both the default and larger example, plus an optional
operator-selected file. Native evidence for those profiles is listed in the resource-policy reference.
These are focused capacity and ownership controls; sustained concurrent pressure and
essential-service fairness remain acceptance work under SEC-07/18.

\section{Extending configuration toward boot resources}

Boot-time installation should reuse this typed policy and its validation, but
must also define a complete authenticated input, reject unknown or incomplete
settings, and establish authority before the first relevant allocation.
Kernel-heap sizing and statically initialized counters currently precede a
general resource-policy installation point. A future implementation must order
those initializations, validate discovered-memory and structural constraints,
then publish one immutable policy before schedulers and services consume it.

Do not implement this by resizing live pools independently. Any future update
mechanism would need a separate account-generation and original-charge contract.
Until that work is implemented and tested, changing capacity means building a
new kernel with a reviewed profile and booting it.

\endinput
